% paper100-copilot-api-bridge.tex — CopilotAPIBridge
%   The API bridge connecting Copilot channels to the JuGeo API layer.
%   Paper 100 of the Judgment Geometry series.
% Compile: pdflatex paper100-copilot-api-bridge
\documentclass[11pt]{article}
\usepackage{lmodern}
% jugeo-common.tex — shared preamble for all Judgment Geometry papers
% Include via: % jugeo-common.tex — shared preamble for all Judgment Geometry papers
% Include via: % jugeo-common.tex — shared preamble for all Judgment Geometry papers
% Include via: \input{jugeo-common}

% ── Packages ────────────────────────────────────────────────────────────
\usepackage{amsmath,amssymb,amsthm,mathtools}
\usepackage{tikz-cd}
\usepackage{listings}
\usepackage{xcolor}
\usepackage{xspace}
\usepackage{booktabs}
\usepackage{multirow}
\usepackage{enumitem}
\usepackage[expansion=false]{microtype}
\usepackage{hyperref}
\usepackage{cleveref}
% \usepackage{thmtools}  % disabled: not available in this TeX installation
\usepackage{float}
\usepackage{caption}
\usepackage{subcaption}
\usepackage{tabularx}
\usepackage{stmaryrd}       % \llbracket, \rrbracket
\usepackage{bbm}            % \mathbbm{1}

% ── Page geometry ───────────────────────────────────────────────────────
\usepackage[margin=1in]{geometry}

% ── Theorem environments ────────────────────────────────────────────────
\theoremstyle{plain}
\newtheorem{theorem}{Theorem}[section]
\newtheorem{lemma}[theorem]{Lemma}
\newtheorem{proposition}[theorem]{Proposition}
\newtheorem{corollary}[theorem]{Corollary}

\theoremstyle{definition}
\newtheorem{definition}[theorem]{Definition}
\newtheorem{example}[theorem]{Example}
\newtheorem{construction}[theorem]{Construction}

\theoremstyle{remark}
\newtheorem{remark}[theorem]{Remark}
\newtheorem{notation}[theorem]{Notation}

% ── Python listings style ──────────────────────────────────────────────
\definecolor{jugeoBlue}{HTML}{2563EB}
\definecolor{jugeoGreen}{HTML}{059669}
\definecolor{jugeoRed}{HTML}{DC2626}
\definecolor{jugeoPurple}{HTML}{7C3AED}
\definecolor{jugeoGray}{HTML}{6B7280}
\definecolor{jugeoBg}{HTML}{F8FAFC}

\lstdefinestyle{jugeo-python}{
  language=Python,
  basicstyle=\ttfamily\small,
  keywordstyle=\color{jugeoBlue}\bfseries,
  stringstyle=\color{jugeoGreen},
  commentstyle=\color{jugeoGray}\itshape,
  numberstyle=\tiny\color{jugeoGray},
  backgroundcolor=\color{jugeoBg},
  numbers=left,
  numbersep=6pt,
  frame=single,
  framerule=0.4pt,
  rulecolor=\color{jugeoGray!40},
  breaklines=true,
  breakatwhitespace=true,
  showstringspaces=false,
  tabsize=4,
  xleftmargin=1.5em,
  framexleftmargin=1.5em,
  morekeywords={dataclass,frozen,field,Enum,IntEnum,auto,Any,
                Mapping,Sequence,Optional,match,case},
  literate={→}{{$\to$}}1 {←}{{$\leftarrow$}}1
           {≤}{{$\leq$}}1 {≥}{{$\geq$}}1 {≠}{{$\neq$}}1
           {∈}{{$\in$}}1 {∀}{{$\forall$}}1 {∃}{{$\exists$}}1
           {⊕}{{$\oplus$}}1 {⊖}{{$\ominus$}}1,
  escapeinside={(*@}{@*)},
}
\lstset{style=jugeo-python}

\lstdefinestyle{jugeo-output}{
  basicstyle=\ttfamily\small,
  backgroundcolor=\color{jugeoBg},
  frame=single,
  framerule=0.4pt,
  rulecolor=\color{jugeoGray!40},
  breaklines=true,
  showstringspaces=false,
  xleftmargin=1.5em,
  framexleftmargin=0.5em,
  numbers=none,
}

% ── JuGeo-specific macros ──────────────────────────────────────────────

% Judgment 8-tuple
\newcommand{\J}{\mathcal{J}}
\newcommand{\Jud}[8]{(#1,\,#2,\,#3,\,#4,\,#5,\,#6,\,#7,\,#8)}
\newcommand{\JudShort}[1]{\J_{#1}}

% Components
\newcommand{\coord}{\mathsf{c}}            % coordinate
\newcommand{\prop}{\varphi}                 % proposition
\newcommand{\carrier}{\mathcal{A}}          % carrier type
\newcommand{\evid}{\mathcal{E}}             % evidence bundle
\newcommand{\oblig}{\mathcal{O}}            % obligations
\newcommand{\obstr}{\mathcal{B}}            % obstructions
\newcommand{\trust}{\mathcal{T}}            % trust annotation
\newcommand{\prov}{\Pi}                     % provenance

% Trust algebra
\newcommand{\Talg}{\mathfrak{T}}            % trust ordered algebra
\newcommand{\Eadm}{\mathcal{E}_{\mathrm{adm}}}
\newcommand{\tleq}{\preceq}                 % trust ordering
\newcommand{\tjoin}{\oplus}                 % conservative join
\newcommand{\tatten}{\ominus}               % attenuation
\newcommand{\tpromote}{\uparrow_\pi}        % promotion
\newcommand{\tdemote}{\downarrow_\chi}       % demotion

% Trust tiers
\newcommand{\Tcontra}{\ensuremath{\mathsf{CONTRADICTED}}}
\newcommand{\Tunver}{\ensuremath{\mathsf{UNVERIFIED}}}
\newcommand{\Tcopilot}{\ensuremath{\mathsf{COPILOT}}}
\newcommand{\Toracle}{\ensuremath{\mathsf{ORACLE}}}
\newcommand{\Truntime}{\ensuremath{\mathsf{RUNTIME}}}
\newcommand{\Tsolver}{\ensuremath{\mathsf{SOLVER}}}
\newcommand{\Tproof}{\ensuremath{\mathsf{PROOF}}}

% Site and geometry
\newcommand{\Site}{\mathcal{S}}             % semantic site
\newcommand{\Cov}{\mathrm{Cov}}            % covering family
\newcommand{\Ob}{\mathrm{Ob}}              % objects
\newcommand{\Mor}{\mathrm{Mor}}            % morphisms
\newcommand{\res}{\mathrm{res}}            % restriction
\newcommand{\Cech}{\check{C}}              % Čech complex
\newcommand{\Hcoh}[1]{H^{#1}}             % cohomology group

% Descent
\newcommand{\descent}{\mathrm{desc}}
\newcommand{\glue}{\mathrm{glue}}
\newcommand{\overlap}[2]{#1 \cap #2}

% Semantic moves
\newcommand{\VERIFY}{\textsc{Verify}}
\newcommand{\CONSTRUCT}{\textsc{Construct}}
\newcommand{\NORMALIZE}{\textsc{Normalize}}
\newcommand{\GLUE}{\textsc{Glue}}
\newcommand{\RESTRICT}{\textsc{Restrict}}
\newcommand{\TRANSPORT}{\textsc{Transport}}
\newcommand{\REPAIR}{\textsc{Repair}}
\newcommand{\PROMOTE}{\textsc{Promote}}
\newcommand{\CHALLENGE}{\textsc{Challenge}}
\newcommand{\ESCALATE}{\textsc{Escalate}}

% Fragments
\newcommand{\QFLIA}{\texttt{QF\_LIA}}
\newcommand{\QFLRA}{\texttt{QF\_LRA}}
\newcommand{\QFBV}{\texttt{QF\_BV}}
\newcommand{\QFUF}{\texttt{QF\_UF}}

% Misc
\newcommand{\Python}{\textsc{Python}}
\newcommand{\JuGeo}{\textsc{JuGeo}}
\newcommand{\LEAN}{\textsc{Lean}\,4}
\newcommand{\Fstar}{F$^{\star}$}
\newcommand{\Coq}{\textsc{Coq}}
\newcommand{\Dafny}{\textsc{Dafny}}

% Common shortcuts
\newcommand{\ie}{i.e.\@\xspace}
\newcommand{\eg}{e.g.\@\xspace}
\newcommand{\cf}{cf.\@\xspace}
\newcommand{\wrt}{w.r.t.\@\xspace}

% Evaluation shortcuts
\newcommand{\numcases}{300}
\newcommand{\numspec}{100}
\newcommand{\numeq}{100}
\newcommand{\numbug}{100}
\newcommand{\medtime}{6.5\,ms}
\newcommand{\totaltime}{1.949\,s}


% ── Packages ────────────────────────────────────────────────────────────
\usepackage{amsmath,amssymb,amsthm,mathtools}
\usepackage{tikz-cd}
\usepackage{listings}
\usepackage{xcolor}
\usepackage{xspace}
\usepackage{booktabs}
\usepackage{multirow}
\usepackage{enumitem}
\usepackage[expansion=false]{microtype}
\usepackage{hyperref}
\usepackage{cleveref}
% \usepackage{thmtools}  % disabled: not available in this TeX installation
\usepackage{float}
\usepackage{caption}
\usepackage{subcaption}
\usepackage{tabularx}
\usepackage{stmaryrd}       % \llbracket, \rrbracket
\usepackage{bbm}            % \mathbbm{1}

% ── Page geometry ───────────────────────────────────────────────────────
\usepackage[margin=1in]{geometry}

% ── Theorem environments ────────────────────────────────────────────────
\theoremstyle{plain}
\newtheorem{theorem}{Theorem}[section]
\newtheorem{lemma}[theorem]{Lemma}
\newtheorem{proposition}[theorem]{Proposition}
\newtheorem{corollary}[theorem]{Corollary}

\theoremstyle{definition}
\newtheorem{definition}[theorem]{Definition}
\newtheorem{example}[theorem]{Example}
\newtheorem{construction}[theorem]{Construction}

\theoremstyle{remark}
\newtheorem{remark}[theorem]{Remark}
\newtheorem{notation}[theorem]{Notation}

% ── Python listings style ──────────────────────────────────────────────
\definecolor{jugeoBlue}{HTML}{2563EB}
\definecolor{jugeoGreen}{HTML}{059669}
\definecolor{jugeoRed}{HTML}{DC2626}
\definecolor{jugeoPurple}{HTML}{7C3AED}
\definecolor{jugeoGray}{HTML}{6B7280}
\definecolor{jugeoBg}{HTML}{F8FAFC}

\lstdefinestyle{jugeo-python}{
  language=Python,
  basicstyle=\ttfamily\small,
  keywordstyle=\color{jugeoBlue}\bfseries,
  stringstyle=\color{jugeoGreen},
  commentstyle=\color{jugeoGray}\itshape,
  numberstyle=\tiny\color{jugeoGray},
  backgroundcolor=\color{jugeoBg},
  numbers=left,
  numbersep=6pt,
  frame=single,
  framerule=0.4pt,
  rulecolor=\color{jugeoGray!40},
  breaklines=true,
  breakatwhitespace=true,
  showstringspaces=false,
  tabsize=4,
  xleftmargin=1.5em,
  framexleftmargin=1.5em,
  morekeywords={dataclass,frozen,field,Enum,IntEnum,auto,Any,
                Mapping,Sequence,Optional,match,case},
  literate={→}{{$\to$}}1 {←}{{$\leftarrow$}}1
           {≤}{{$\leq$}}1 {≥}{{$\geq$}}1 {≠}{{$\neq$}}1
           {∈}{{$\in$}}1 {∀}{{$\forall$}}1 {∃}{{$\exists$}}1
           {⊕}{{$\oplus$}}1 {⊖}{{$\ominus$}}1,
  escapeinside={(*@}{@*)},
}
\lstset{style=jugeo-python}

\lstdefinestyle{jugeo-output}{
  basicstyle=\ttfamily\small,
  backgroundcolor=\color{jugeoBg},
  frame=single,
  framerule=0.4pt,
  rulecolor=\color{jugeoGray!40},
  breaklines=true,
  showstringspaces=false,
  xleftmargin=1.5em,
  framexleftmargin=0.5em,
  numbers=none,
}

% ── JuGeo-specific macros ──────────────────────────────────────────────

% Judgment 8-tuple
\newcommand{\J}{\mathcal{J}}
\newcommand{\Jud}[8]{(#1,\,#2,\,#3,\,#4,\,#5,\,#6,\,#7,\,#8)}
\newcommand{\JudShort}[1]{\J_{#1}}

% Components
\newcommand{\coord}{\mathsf{c}}            % coordinate
\newcommand{\prop}{\varphi}                 % proposition
\newcommand{\carrier}{\mathcal{A}}          % carrier type
\newcommand{\evid}{\mathcal{E}}             % evidence bundle
\newcommand{\oblig}{\mathcal{O}}            % obligations
\newcommand{\obstr}{\mathcal{B}}            % obstructions
\newcommand{\trust}{\mathcal{T}}            % trust annotation
\newcommand{\prov}{\Pi}                     % provenance

% Trust algebra
\newcommand{\Talg}{\mathfrak{T}}            % trust ordered algebra
\newcommand{\Eadm}{\mathcal{E}_{\mathrm{adm}}}
\newcommand{\tleq}{\preceq}                 % trust ordering
\newcommand{\tjoin}{\oplus}                 % conservative join
\newcommand{\tatten}{\ominus}               % attenuation
\newcommand{\tpromote}{\uparrow_\pi}        % promotion
\newcommand{\tdemote}{\downarrow_\chi}       % demotion

% Trust tiers
\newcommand{\Tcontra}{\ensuremath{\mathsf{CONTRADICTED}}}
\newcommand{\Tunver}{\ensuremath{\mathsf{UNVERIFIED}}}
\newcommand{\Tcopilot}{\ensuremath{\mathsf{COPILOT}}}
\newcommand{\Toracle}{\ensuremath{\mathsf{ORACLE}}}
\newcommand{\Truntime}{\ensuremath{\mathsf{RUNTIME}}}
\newcommand{\Tsolver}{\ensuremath{\mathsf{SOLVER}}}
\newcommand{\Tproof}{\ensuremath{\mathsf{PROOF}}}

% Site and geometry
\newcommand{\Site}{\mathcal{S}}             % semantic site
\newcommand{\Cov}{\mathrm{Cov}}            % covering family
\newcommand{\Ob}{\mathrm{Ob}}              % objects
\newcommand{\Mor}{\mathrm{Mor}}            % morphisms
\newcommand{\res}{\mathrm{res}}            % restriction
\newcommand{\Cech}{\check{C}}              % Čech complex
\newcommand{\Hcoh}[1]{H^{#1}}             % cohomology group

% Descent
\newcommand{\descent}{\mathrm{desc}}
\newcommand{\glue}{\mathrm{glue}}
\newcommand{\overlap}[2]{#1 \cap #2}

% Semantic moves
\newcommand{\VERIFY}{\textsc{Verify}}
\newcommand{\CONSTRUCT}{\textsc{Construct}}
\newcommand{\NORMALIZE}{\textsc{Normalize}}
\newcommand{\GLUE}{\textsc{Glue}}
\newcommand{\RESTRICT}{\textsc{Restrict}}
\newcommand{\TRANSPORT}{\textsc{Transport}}
\newcommand{\REPAIR}{\textsc{Repair}}
\newcommand{\PROMOTE}{\textsc{Promote}}
\newcommand{\CHALLENGE}{\textsc{Challenge}}
\newcommand{\ESCALATE}{\textsc{Escalate}}

% Fragments
\newcommand{\QFLIA}{\texttt{QF\_LIA}}
\newcommand{\QFLRA}{\texttt{QF\_LRA}}
\newcommand{\QFBV}{\texttt{QF\_BV}}
\newcommand{\QFUF}{\texttt{QF\_UF}}

% Misc
\newcommand{\Python}{\textsc{Python}}
\newcommand{\JuGeo}{\textsc{JuGeo}}
\newcommand{\LEAN}{\textsc{Lean}\,4}
\newcommand{\Fstar}{F$^{\star}$}
\newcommand{\Coq}{\textsc{Coq}}
\newcommand{\Dafny}{\textsc{Dafny}}

% Common shortcuts
\newcommand{\ie}{i.e.\@\xspace}
\newcommand{\eg}{e.g.\@\xspace}
\newcommand{\cf}{cf.\@\xspace}
\newcommand{\wrt}{w.r.t.\@\xspace}

% Evaluation shortcuts
\newcommand{\numcases}{300}
\newcommand{\numspec}{100}
\newcommand{\numeq}{100}
\newcommand{\numbug}{100}
\newcommand{\medtime}{6.5\,ms}
\newcommand{\totaltime}{1.949\,s}


% ── Packages ────────────────────────────────────────────────────────────
\usepackage{amsmath,amssymb,amsthm,mathtools}
\usepackage{tikz-cd}
\usepackage{listings}
\usepackage{xcolor}
\usepackage{xspace}
\usepackage{booktabs}
\usepackage{multirow}
\usepackage{enumitem}
\usepackage[expansion=false]{microtype}
\usepackage{hyperref}
\usepackage{cleveref}
% \usepackage{thmtools}  % disabled: not available in this TeX installation
\usepackage{float}
\usepackage{caption}
\usepackage{subcaption}
\usepackage{tabularx}
\usepackage{stmaryrd}       % \llbracket, \rrbracket
\usepackage{bbm}            % \mathbbm{1}

% ── Page geometry ───────────────────────────────────────────────────────
\usepackage[margin=1in]{geometry}

% ── Theorem environments ────────────────────────────────────────────────
\theoremstyle{plain}
\newtheorem{theorem}{Theorem}[section]
\newtheorem{lemma}[theorem]{Lemma}
\newtheorem{proposition}[theorem]{Proposition}
\newtheorem{corollary}[theorem]{Corollary}

\theoremstyle{definition}
\newtheorem{definition}[theorem]{Definition}
\newtheorem{example}[theorem]{Example}
\newtheorem{construction}[theorem]{Construction}

\theoremstyle{remark}
\newtheorem{remark}[theorem]{Remark}
\newtheorem{notation}[theorem]{Notation}

% ── Python listings style ──────────────────────────────────────────────
\definecolor{jugeoBlue}{HTML}{2563EB}
\definecolor{jugeoGreen}{HTML}{059669}
\definecolor{jugeoRed}{HTML}{DC2626}
\definecolor{jugeoPurple}{HTML}{7C3AED}
\definecolor{jugeoGray}{HTML}{6B7280}
\definecolor{jugeoBg}{HTML}{F8FAFC}

\lstdefinestyle{jugeo-python}{
  language=Python,
  basicstyle=\ttfamily\small,
  keywordstyle=\color{jugeoBlue}\bfseries,
  stringstyle=\color{jugeoGreen},
  commentstyle=\color{jugeoGray}\itshape,
  numberstyle=\tiny\color{jugeoGray},
  backgroundcolor=\color{jugeoBg},
  numbers=left,
  numbersep=6pt,
  frame=single,
  framerule=0.4pt,
  rulecolor=\color{jugeoGray!40},
  breaklines=true,
  breakatwhitespace=true,
  showstringspaces=false,
  tabsize=4,
  xleftmargin=1.5em,
  framexleftmargin=1.5em,
  morekeywords={dataclass,frozen,field,Enum,IntEnum,auto,Any,
                Mapping,Sequence,Optional,match,case},
  literate={→}{{$\to$}}1 {←}{{$\leftarrow$}}1
           {≤}{{$\leq$}}1 {≥}{{$\geq$}}1 {≠}{{$\neq$}}1
           {∈}{{$\in$}}1 {∀}{{$\forall$}}1 {∃}{{$\exists$}}1
           {⊕}{{$\oplus$}}1 {⊖}{{$\ominus$}}1,
  escapeinside={(*@}{@*)},
}
\lstset{style=jugeo-python}

\lstdefinestyle{jugeo-output}{
  basicstyle=\ttfamily\small,
  backgroundcolor=\color{jugeoBg},
  frame=single,
  framerule=0.4pt,
  rulecolor=\color{jugeoGray!40},
  breaklines=true,
  showstringspaces=false,
  xleftmargin=1.5em,
  framexleftmargin=0.5em,
  numbers=none,
}

% ── JuGeo-specific macros ──────────────────────────────────────────────

% Judgment 8-tuple
\newcommand{\J}{\mathcal{J}}
\newcommand{\Jud}[8]{(#1,\,#2,\,#3,\,#4,\,#5,\,#6,\,#7,\,#8)}
\newcommand{\JudShort}[1]{\J_{#1}}

% Components
\newcommand{\coord}{\mathsf{c}}            % coordinate
\newcommand{\prop}{\varphi}                 % proposition
\newcommand{\carrier}{\mathcal{A}}          % carrier type
\newcommand{\evid}{\mathcal{E}}             % evidence bundle
\newcommand{\oblig}{\mathcal{O}}            % obligations
\newcommand{\obstr}{\mathcal{B}}            % obstructions
\newcommand{\trust}{\mathcal{T}}            % trust annotation
\newcommand{\prov}{\Pi}                     % provenance

% Trust algebra
\newcommand{\Talg}{\mathfrak{T}}            % trust ordered algebra
\newcommand{\Eadm}{\mathcal{E}_{\mathrm{adm}}}
\newcommand{\tleq}{\preceq}                 % trust ordering
\newcommand{\tjoin}{\oplus}                 % conservative join
\newcommand{\tatten}{\ominus}               % attenuation
\newcommand{\tpromote}{\uparrow_\pi}        % promotion
\newcommand{\tdemote}{\downarrow_\chi}       % demotion

% Trust tiers
\newcommand{\Tcontra}{\ensuremath{\mathsf{CONTRADICTED}}}
\newcommand{\Tunver}{\ensuremath{\mathsf{UNVERIFIED}}}
\newcommand{\Tcopilot}{\ensuremath{\mathsf{COPILOT}}}
\newcommand{\Toracle}{\ensuremath{\mathsf{ORACLE}}}
\newcommand{\Truntime}{\ensuremath{\mathsf{RUNTIME}}}
\newcommand{\Tsolver}{\ensuremath{\mathsf{SOLVER}}}
\newcommand{\Tproof}{\ensuremath{\mathsf{PROOF}}}

% Site and geometry
\newcommand{\Site}{\mathcal{S}}             % semantic site
\newcommand{\Cov}{\mathrm{Cov}}            % covering family
\newcommand{\Ob}{\mathrm{Ob}}              % objects
\newcommand{\Mor}{\mathrm{Mor}}            % morphisms
\newcommand{\res}{\mathrm{res}}            % restriction
\newcommand{\Cech}{\check{C}}              % Čech complex
\newcommand{\Hcoh}[1]{H^{#1}}             % cohomology group

% Descent
\newcommand{\descent}{\mathrm{desc}}
\newcommand{\glue}{\mathrm{glue}}
\newcommand{\overlap}[2]{#1 \cap #2}

% Semantic moves
\newcommand{\VERIFY}{\textsc{Verify}}
\newcommand{\CONSTRUCT}{\textsc{Construct}}
\newcommand{\NORMALIZE}{\textsc{Normalize}}
\newcommand{\GLUE}{\textsc{Glue}}
\newcommand{\RESTRICT}{\textsc{Restrict}}
\newcommand{\TRANSPORT}{\textsc{Transport}}
\newcommand{\REPAIR}{\textsc{Repair}}
\newcommand{\PROMOTE}{\textsc{Promote}}
\newcommand{\CHALLENGE}{\textsc{Challenge}}
\newcommand{\ESCALATE}{\textsc{Escalate}}

% Fragments
\newcommand{\QFLIA}{\texttt{QF\_LIA}}
\newcommand{\QFLRA}{\texttt{QF\_LRA}}
\newcommand{\QFBV}{\texttt{QF\_BV}}
\newcommand{\QFUF}{\texttt{QF\_UF}}

% Misc
\newcommand{\Python}{\textsc{Python}}
\newcommand{\JuGeo}{\textsc{JuGeo}}
\newcommand{\LEAN}{\textsc{Lean}\,4}
\newcommand{\Fstar}{F$^{\star}$}
\newcommand{\Coq}{\textsc{Coq}}
\newcommand{\Dafny}{\textsc{Dafny}}

% Common shortcuts
\newcommand{\ie}{i.e.\@\xspace}
\newcommand{\eg}{e.g.\@\xspace}
\newcommand{\cf}{cf.\@\xspace}
\newcommand{\wrt}{w.r.t.\@\xspace}

% Evaluation shortcuts
\newcommand{\numcases}{300}
\newcommand{\numspec}{100}
\newcommand{\numeq}{100}
\newcommand{\numbug}{100}
\newcommand{\medtime}{6.5\,ms}
\newcommand{\totaltime}{1.949\,s}


% ——— Paper-specific macros (letters only) ———
\newcommand{\Bridge}{\mathsf{CopilotAPIBridge}}
\newcommand{\JuGeoAPI}{\mathsf{JuGeoAPI}}
\newcommand{\APIReq}{\mathsf{APIRequest}}
\newcommand{\APIResp}{\mathsf{APIResponse}}
\newcommand{\APISess}{\mathsf{APISession}}
\newcommand{\APILog}{\mathsf{APIEventLog}}
\newcommand{\APIAuth}{\mathsf{APIAuthenticator}}
\newcommand{\APIRate}{\mathsf{APIRateLimiter}}
\newcommand{\APIVal}{\mathsf{APIValidator}}
\newcommand{\OpKind}{\mathsf{OperationKind}}
\newcommand{\ReqStat}{\mathsf{RequestStatus}}
\newcommand{\CopChan}{\mathsf{CopilotChannel}}
\newcommand{\ChanCfg}{\mathsf{ChannelConfiguration}}
\newcommand{\ChanPool}{\mathsf{ChannelPool}}
\newcommand{\ChanMon}{\mathsf{ChannelMonitor}}
\newcommand{\ChanFed}{\mathsf{ChannelFederation}}
\newcommand{\EvReq}{\mathsf{EvidenceRequest}}
\newcommand{\EvResp}{\mathsf{EvidenceResponse}}
\newcommand{\EvRec}{\mathsf{EvidenceRecord}}
\newcommand{\SuppRoute}{\mathsf{SupportRoute}}

\title{\textbf{CopilotAPIBridge:\\
Bridging Copilot Channels and the JuGeo API Layer}}
\author{Halley Young}
\date{Paper~100 --- Judgment Geometry Series}

\begin{document}
\maketitle

% ============================================================
%  Abstract
% ============================================================
\begin{abstract}
The $\Bridge$ is the central integration point between the
Copilot evidence channel and the \JuGeo{} API layer.  It
translates API-level operations (verify, construct, descend,
federate, inspect) into Copilot channel queries, enforces
the $\Tcopilot$ trust ceiling, optionally requires
corroboration from higher-trust channels, and logs all
interactions.  This paper provides a comprehensive formal
treatment: we model the bridge as a trust-bounded natural
transformation between the API functor and the channel
functor, prove five theorems (trust ceiling enforcement,
corroboration soundness, event log completeness, session
isolation, and rate-limit compliance), and present four
\Python{} listings covering the full bridge lifecycle.
An appendix collects auxiliary lemmas and proofs.
\end{abstract}

% ============================================================
%  Section 1 — Introduction
% ============================================================
\section{Introduction}\label{sec:intro}

The \JuGeo{} framework exposes a public API through the
$\JuGeoAPI$ class.  This API supports five operation kinds:
$\OpKind.\texttt{VERIFY}$, $\OpKind.\texttt{CONSTRUCT}$,
$\OpKind.\texttt{DESCEND}$, $\OpKind.\texttt{FEDERATE}$,
and $\OpKind.\texttt{INSPECT}$.  Internally, many of these
operations benefit from Copilot assistance: the language
model can suggest verification strategies, construction
hints, descent directions, and inspection summaries.

The $\Bridge$ mediates this integration.  It receives
API-level requests, formulates appropriate prompts, queries
the Copilot channel, post-processes the responses (clamping
trust, checking corroboration), and returns API-level
responses.  All interactions are logged in the $\APILog$.

\paragraph{Contributions.}
\begin{enumerate}
\item A categorical model of the bridge as a natural
      transformation between two functors
      (\S\ref{sec:categorical}).
\item A detailed description of the bridge's trust
      enforcement and corroboration mechanisms
      (\S\ref{sec:trustenforce}).
\item The design of the event logging and session management
      protocols (\S\ref{sec:logging}).
\item Five theorems: trust ceiling enforcement, corroboration
      soundness, event log completeness, session isolation,
      and rate-limit compliance (\S\ref{sec:formal}).
\item Four \Python{} listings showing the full bridge
      implementation (\S\ref{sec:code}).
\end{enumerate}

\paragraph{Paper outline.}
Section~\ref{sec:background} recalls the API and channel layers.
Section~\ref{sec:categorical} presents the categorical model.
Section~\ref{sec:trustenforce} describes trust enforcement.
Section~\ref{sec:corroboration} treats corroboration.
Section~\ref{sec:logging} covers event logging.
Section~\ref{sec:code} gives the implementation.
Section~\ref{sec:formal} states and proves the main theorems.
Section~\ref{sec:discussion} discusses limitations.
Section~\ref{sec:related} surveys related work.
Section~\ref{sec:conclusion} concludes.
Appendix~\ref{app:proofs} collects auxiliary proofs.

% ============================================================
%  Section 2 — Background
% ============================================================
\section{Background}\label{sec:background}

\subsection{The JuGeo API layer}

The API layer consists of:
\begin{itemize}
\item $\JuGeoAPI$: the main entry point, supporting
      \texttt{verify}, \texttt{construct}, \texttt{descend},
      \texttt{federate}, \texttt{inspect}, and
      \texttt{copilot\_query}.
\item $\APIReq$: encapsulates a request with id, operation,
      coordinate, proposition, budget, deadline, and metadata.
\item $\APIResp$: encapsulates a response with status,
      evidence, trust, and residuals.
\item $\APISess$: manages a sequence of related requests.
\item $\APILog$: records all API events for auditing.
\item $\APIAuth$: authenticates callers.
\item $\APIRate$: enforces rate limits.
\item $\APIVal$: validates request parameters.
\end{itemize}

\begin{definition}[Operation kind]\label{def:opkind}
$\OpKind$ is an enumeration:
$\{\texttt{VERIFY}, \texttt{CONSTRUCT}, \texttt{DESCEND},
\texttt{FEDERATE}, \texttt{INSPECT}\}$.
Each kind determines the prompt template used by the bridge.
\end{definition}

\subsection{The evidence channel layer}

The channel layer (Paper~88) provides:
\begin{itemize}
\item $\CopChan$: queries a language model with trust ceiling
      $\Tcopilot$.
\item $\ChanPool$: manages registered channels.
\item $\ChanMon$: records request--response pairs.
\item $\ChanFed$: collects evidence from multiple channels.
\end{itemize}

\subsection{Trust lattice}

\[
  \Tcontra \;\tleq\; \Tunver \;\tleq\; \Tcopilot
  \;\tleq\; \Toracle \;\tleq\; \Truntime
  \;\tleq\; \Tsolver \;\tleq\; \Tproof.
\]
The bridge operates at $\Tcopilot$ and never produces
evidence at a higher trust tier.

% ============================================================
%  Section 3 — Categorical Model
% ============================================================
\section{Categorical Model}\label{sec:categorical}

\subsection{The API functor}

\begin{definition}[API category]\label{def:apicat}
Let $\mathbf{API}$ be the category whose objects are
$\OpKind$ values and whose morphisms are operation
refinements (\eg{} $\texttt{VERIFY} \to \texttt{INSPECT}$
when verification includes inspection as a sub-operation).
\end{definition}

\begin{definition}[API functor]\label{def:apifunctor}
The \emph{API functor} $F_A \colon \mathbf{API} \to \mathbf{Set}$
sends each operation kind~$o$ to the set of valid $\APIReq$
values with $\texttt{operation} = o$, and each morphism
$o \to o'$ to the function that restricts a request of
kind~$o$ to a sub-request of kind~$o'$.
\end{definition}

\subsection{The channel functor}

\begin{definition}[Channel functor]\label{def:chanfunctor}
The \emph{channel functor}
$F_C \colon \mathbf{API} \to \mathbf{Set}$ sends each
operation kind~$o$ to the set of $\EvReq$ values that
the bridge generates for operation~$o$, and each morphism
to the corresponding restriction on evidence requests.
\end{definition}

\subsection{The bridge as a natural transformation}

\begin{definition}[Bridge natural transformation]%
\label{def:bridgenat}
The $\Bridge$ is a natural transformation
$\eta \colon F_A \Rightarrow F_C$ where for each
operation kind~$o$, $\eta_o$ maps an $\APIReq$ to an
$\EvReq$.  Naturality means: for every morphism
$o \to o'$ in $\mathbf{API}$, the following diagram
commutes:
\[
\begin{array}{ccc}
  F_A(o) & \xrightarrow{\eta_o} & F_C(o) \\
  \downarrow & & \downarrow \\
  F_A(o') & \xrightarrow{\eta_{o'}} & F_C(o')
\end{array}
\]
\end{definition}

\begin{remark}
Naturality captures the intuition that the bridge's
translation from API requests to channel requests is
coherent: restricting an API request and then translating
gives the same result as translating and then restricting
the evidence request.
\end{remark}

% ============================================================
%  Section 4 — Trust Enforcement
% ============================================================
\section{Trust Enforcement}\label{sec:trustenforce}

\subsection{Trust ceiling mechanism}

\begin{definition}[Trust enforcement]\label{def:trustenforce}
The bridge enforces the trust ceiling by calling
$\CopChan.\texttt{apply\_trust\_ceiling}()$ on every
response.  This caps the trust of the evidence at
$\Tcopilot$.  Additionally, the bridge calls
$\texttt{enforce\_trust\_ceiling}()$ on the bridge
object itself, which verifies that no response in
the current session exceeds $\Tcopilot$.
\end{definition}

\begin{lemma}[Bridge trust ceiling]\label{lem:bridgeceil}
For every API request processed by the bridge, the
corresponding $\APIResp$ has trust $\tleq \Tcopilot$.
\end{lemma}
\begin{proof}
The bridge calls \texttt{apply\_trust\_ceiling} on the
Copilot channel response, which caps trust at $\Tcopilot$.
The bridge then wraps this in an $\APIResp$ without
escalating trust.  Therefore the API response trust is
at most $\Tcopilot$.
\end{proof}

\subsection{Trust clamping}

\begin{definition}[Trust clamping]\label{def:trustclamp}
Given an $\EvResp$ with trust~$\tau$, the clamped response
has trust $\min(\tau, \Tcopilot)$.  The method
$\EvResp.\texttt{clamp\_trust}(\texttt{"proposal"})$
performs this operation.
\end{definition}

\begin{lemma}[Clamping is idempotent]\label{lem:clampidemp}
Clamping a response that has already been clamped yields
the same response: $\mathrm{clamp}(\mathrm{clamp}(r)) =
\mathrm{clamp}(r)$.
\end{lemma}
\begin{proof}
After the first clamp, $\mathrm{trust}(r) \leq \Tcopilot$.
The second clamp computes $\min(\mathrm{trust}(r), \Tcopilot)
= \mathrm{trust}(r)$.  Therefore the response is unchanged.
\end{proof}

% ============================================================
%  Section 5 — Corroboration
% ============================================================
\section{Corroboration}\label{sec:corroboration}

\subsection{Corroboration protocol}

When the bridge is configured with
$\texttt{require\_corroboration} = \texttt{True}$, every
Copilot response must be corroborated by evidence from a
higher-trust channel before being returned to the API layer.

\begin{definition}[Corroboration]\label{def:corroboration}
A Copilot response~$r$ at trust $\Tcopilot$ is
\emph{corroborated} if there exists a response~$r'$ from
a channel with trust $\tau' > \Tcopilot$ such that
$r$ and $r'$ are consistent (\ie{} they do not contradict
each other on the same proposition).
\end{definition}

\begin{definition}[Consistency]\label{def:consistency}
Two evidence responses $r$ and $r'$ are \emph{consistent}
if they do not assign contradictory polarities to the same
proposition.  Formally:
\[
  \forall \varphi.\; r(\varphi) \neq \bot \wedge
  r'(\varphi) \neq \bot \implies
  r(\varphi) = r'(\varphi).
\]
\end{definition}

\begin{lemma}[Corroboration increases effective trust]%
\label{lem:corrobtrust}
If a Copilot response at $\Tcopilot$ is corroborated by
a solver response at $\Tsolver$, the effective trust of
the corroborated evidence is $\Tsolver$ for the
corroborated propositions.
\end{lemma}
\begin{proof}
The corroborated propositions have independent evidence at
$\Tsolver$.  The Copilot response serves as a hint that
guided the solver query; the trust of the final evidence
is determined by the solver, which is at $\Tsolver$.
\end{proof}

\subsection{Corroboration failure}

If corroboration fails (no higher-trust channel confirms
the Copilot response), the bridge may:
\begin{enumerate}
\item Return the uncorroborated response at $\Tcopilot$, or
\item Return an empty response indicating that the Copilot
      suggestion could not be corroborated, or
\item Log the failure and continue.
\end{enumerate}
The choice depends on the bridge configuration and the
operation kind.

% ============================================================
%  Section 6 — Event Logging
% ============================================================
\section{Event Logging and Session Management}\label{sec:logging}

\subsection{Event log}

The $\APILog$ records every API event: request submission,
Copilot query, channel response, trust clamping, corroboration
attempt, and final API response.

\begin{definition}[API event]\label{def:apievent}
An \emph{API event} is a tuple
$(t, k, d)$ where $t$ is a timestamp, $k$ is the event
kind (request, query, response, clamp, corroborate, error),
and $d$ is the event data.
\end{definition}

\begin{lemma}[Event log completeness]\label{lem:logcomp}
Every bridge invocation generates at least three events:
request received, Copilot query dispatched, and response
returned.
\end{lemma}
\begin{proof}
The bridge's \texttt{query} method logs the incoming request,
calls the Copilot channel (logging the query), and logs the
response before returning.  Even if the query fails, an error
event is logged.
\end{proof}

\subsection{Session management}

The $\APISess$ tracks a sequence of related API requests.
Each session has an id, an open/closed state, and a
checkpoint mechanism.

\begin{definition}[Session lifecycle]\label{def:sesslife}
A session transitions through states:
$\mathsf{open} \to \mathsf{checkpointed} \to \mathsf{closed}$.
The \texttt{register\_request} method adds requests.
The \texttt{checkpoint} method serialises the session state.
The \texttt{close} method finalises the session.
\end{definition}

\begin{lemma}[Session isolation]\label{lem:sessiso}
Distinct sessions do not share mutable state.
\end{lemma}
\begin{proof}
Each $\APISess$ instance maintains a private request list,
checkpoint state, and session id.  No global mutable state
is shared between sessions.
\end{proof}

% ============================================================
%  Section 7 — Implementation
% ============================================================
\section{Implementation}\label{sec:code}

\subsection{Core bridge implementation}

The following listing shows the complete $\Bridge$
class with query, trust enforcement, and corroboration.

\begin{lstlisting}[style=jugeo-python,
  caption={Core CopilotAPIBridge implementation},
  label={lst:bridge}]
from jugeo.interfaces.api import (
    CopilotAPIBridge,
    APIEventLog,
    APIRequest,
    APIResponse,
    OperationKind,
    RequestStatus,
    APISession,
    JuGeoAPI,
    APIValidator,
    APIRateLimiter,
    APIAuthenticator,
)
from jugeo.evidence.channels import (
    CopilotChannel,
    ChannelConfiguration,
    EvidenceChannel,
    EvidenceRequest,
    EvidenceResponse,
    ChannelMonitor,
    SolverChannel,
    ChannelPool,
)

class ExtendedCopilotAPIBridge:
    """Extended bridge with corroboration and session management."""

    def __init__(self, channel: CopilotChannel, event_log: APIEventLog,
                 require_corroboration=False, pool=None):
        self.bridge = CopilotAPIBridge(
            channel=channel,
            event_log=event_log,
            require_corroboration=require_corroboration,
        )
        self.channel = channel
        self.event_log = event_log
        self.require_corroboration = require_corroboration
        self.pool = pool

    def query_for_operation(self, operation, coordinate, proposition):
        prompt = self._build_prompt(operation, coordinate, proposition)
        result = self.bridge.query(prompt)
        self.bridge.enforce_trust_ceiling()
        self.event_log.record({
            "event": "bridge_query",
            "operation": str(operation),
            "coordinate": coordinate,
        })
        if self.require_corroboration and self.pool:
            corroborated = self._corroborate(result, coordinate)
            if corroborated:
                return corroborated
        return result

    def verify_via_bridge(self, api: JuGeoAPI, coordinate, proposition):
        suggestion = self.query_for_operation(
            OperationKind.VERIFY, coordinate, proposition
        )
        result = api.verify(coordinate=coordinate, proposition=proposition)
        self.event_log.record({
            "event": "verify_with_suggestion",
            "suggestion": str(suggestion),
            "result": str(result),
        })
        return result

    def construct_via_bridge(self, api: JuGeoAPI, coordinate, proposition):
        suggestion = self.query_for_operation(
            OperationKind.CONSTRUCT, coordinate, proposition
        )
        result = api.construct(coordinate=coordinate, proposition=proposition)
        self.event_log.record({
            "event": "construct_with_suggestion",
            "result": str(result),
        })
        return result

    def inspect_via_bridge(self, api: JuGeoAPI, coordinate):
        suggestion = self.query_for_operation(
            OperationKind.INSPECT, coordinate, None
        )
        result = api.inspect(coordinate=coordinate)
        return {"suggestion": suggestion, "inspection": result}

    def get_stats(self):
        return self.bridge.stats()

    def _build_prompt(self, operation, coordinate, proposition):
        return (
            f"Operation: {operation}, "
            f"Coordinate: {coordinate}, "
            f"Proposition: {proposition}"
        )

    def _corroborate(self, copilot_result, coordinate):
        solver = self.pool.get_channel("solver") if self.pool else None
        if solver is None:
            return None
        solver_req = EvidenceRequest(
            request_id=f"corroborate-{coordinate}",
            coordinate=coordinate,
            proposition=str(copilot_result),
            required_kind=None,
            proposition_kind="corroboration",
            preferred_channel=EvidenceChannel.SOLVER,
            fallback_channels=[],
            deadline_ms=10000,
            budget=20000,
            metadata={"source": "corroboration"},
        )
        solver_resp = solver.handle_request(solver_req)
        if solver_resp and not solver_resp.is_empty():
            self.event_log.record({
                "event": "corroboration_success",
                "coordinate": coordinate,
            })
            return solver_resp
        self.event_log.record({
            "event": "corroboration_failed",
            "coordinate": coordinate,
        })
        return None
\end{lstlisting}

\subsection{Full API integration}

The next listing shows how $\JuGeoAPI$ is instantiated
with the bridge and supporting components.

\begin{lstlisting}[style=jugeo-python,
  caption={Full JuGeoAPI instantiation with CopilotAPIBridge},
  label={lst:apisetup}]
from jugeo.interfaces.api import (
    JuGeoAPI,
    CopilotAPIBridge,
    APIEventLog,
    APIAuthenticator,
    APIRateLimiter,
    APIValidator,
    APISession,
    OperationKind,
    APIRequest,
    APIResponse,
)
from jugeo.evidence.channels import (
    CopilotChannel,
    SolverChannel,
    RuntimeChannel,
    ChannelPool,
    ChannelRouter,
    ChannelMonitor,
    ChannelConfiguration,
    EvidenceChannel,
)

def create_full_api(config_copilot, config_solver):
    """Create a fully configured JuGeoAPI with Copilot bridge."""
    pool = ChannelPool()
    monitor = ChannelMonitor()
    router = ChannelRouter()

    copilot = CopilotChannel(config=config_copilot, model_name="default")
    solver = SolverChannel(config=config_solver, session_id="main")
    runtime = RuntimeChannel()

    pool.register_channel("copilot", copilot)
    pool.register_channel("solver", solver)
    pool.register_channel("runtime", runtime)

    event_log = APIEventLog()
    authenticator = APIAuthenticator()
    authenticator.register_caller("default-user")
    rate_limiter = APIRateLimiter(capacity=100, refill_rate=10)
    validator = APIValidator(max_budget=100000)

    bridge = CopilotAPIBridge(
        channel=copilot,
        event_log=event_log,
        require_corroboration=True,
    )

    api = JuGeoAPI(
        manifest=None,
        authenticator=authenticator,
        rate_limiter=rate_limiter,
        validator=validator,
        router=router,
        event_log=event_log,
        copilot_bridge=bridge,
    )
    return api, pool, monitor

def run_verification_session(api: JuGeoAPI, coordinate, proposition):
    """Run a complete verification session with Copilot assistance."""
    session = api.open_session("verification-session")
    request = APIRequest(
        request_id="ver-001",
        operation=OperationKind.VERIFY,
        coordinate=coordinate,
        proposition=proposition,
        budget=50000,
        deadline=30000,
        metadata={"session": "main"},
    )
    session.register_request(request)

    copilot_hint = api.copilot_query(
        f"Strategy for verifying {proposition} at {coordinate}"
    )
    result = api.verify(coordinate=coordinate, proposition=proposition)

    session.checkpoint()
    diag = api.diagnostics()
    session.close()

    return {
        "result": result,
        "copilot_hint": copilot_hint,
        "diagnostics": diag,
        "session_closed": not session.is_open(),
    }
\end{lstlisting}

\subsection{Federated bridge operation}

The following listing shows how the bridge integrates
with $\ChanFed$ for federated operations.

\begin{lstlisting}[style=jugeo-python,
  caption={Federated bridge operation},
  label={lst:fedbridge}]
from jugeo.evidence.channels import (
    ChannelFederation,
    ChannelPool,
    ChannelRouter,
    ChannelMonitor,
    CopilotChannel,
    SolverChannel,
    EvidenceChannel,
    EvidenceRequest,
    EvidenceResponse,
    EvidenceRecord,
    SupportRoute,
    ChannelSerializer,
)
from jugeo.interfaces.api import (
    CopilotAPIBridge,
    APIEventLog,
    JuGeoAPI,
    OperationKind,
    APIRequest,
)

class FederatedBridgeOperation:
    """Bridge-mediated federated evidence collection."""

    def __init__(self, bridge: CopilotAPIBridge, pool: ChannelPool):
        self.bridge = bridge
        self.pool = pool
        self.federation = ChannelFederation()
        self.router = ChannelRouter()
        self.monitor = ChannelMonitor()

    def federated_query(self, coordinate, proposition):
        copilot_hint = self.bridge.query(
            f"Federate evidence for {proposition} at {coordinate}"
        )
        self.bridge.enforce_trust_ceiling()

        request = EvidenceRequest(
            request_id=f"fed-bridge-{coordinate}",
            coordinate=coordinate,
            proposition=proposition,
            required_kind=None,
            proposition_kind="federated",
            preferred_channel=EvidenceChannel.COPILOT,
            fallback_channels=[
                EvidenceChannel.SOLVER,
                EvidenceChannel.RUNTIME,
            ],
            deadline_ms=20000,
            budget=30000,
            metadata={"bridge_hint": str(copilot_hint)},
        )

        self.monitor.record_request(request)
        best = self.router.find_best_channel(request)
        all_capable = self.router.find_all_capable(request)

        evidence = self.federation.collect_all(request)
        merged = self.federation.merge_by_kind(evidence)
        self.federation.verify_no_escalation(merged)
        composed_trust = self.federation.compose_trust(merged)

        for ev in evidence:
            self.monitor.record_response(ev)

        serialized = ChannelSerializer.request_to_json(request)
        round_trip = ChannelSerializer.round_trip_request(serialized)

        return {
            "merged": merged,
            "trust": composed_trust,
            "hint": copilot_hint,
            "success_rate": self.monitor.success_rate(),
            "latencies": self.monitor.latency_percentiles(),
            "summary": self.monitor.summary(),
        }

    def federate_via_api(self, api: JuGeoAPI, coordinate, proposition):
        result = api.federate(
            coordinate=coordinate,
            proposition=proposition,
        )
        return result
\end{lstlisting}

\subsection{Event log export and diagnostics}

The final listing shows event log management and
diagnostic output.

\begin{lstlisting}[style=jugeo-python,
  caption={Event log export and API diagnostics},
  label={lst:logdiag}]
from jugeo.interfaces.api import (
    JuGeoAPI,
    APIEventLog,
    APISession,
    APIRequest,
    APIResponse,
    OperationKind,
    CopilotAPIBridge,
)
from jugeo.evidence.channels import (
    CopilotChannel,
    ChannelConfiguration,
    EvidenceChannel,
    ChannelMonitor,
)

class BridgeDiagnostics:
    """Diagnostics and auditing for the CopilotAPIBridge."""

    def __init__(self, api: JuGeoAPI, bridge: CopilotAPIBridge,
                 event_log: APIEventLog):
        self.api = api
        self.bridge = bridge
        self.event_log = event_log

    def export_log(self, path=None):
        exported = self.event_log.export()
        if path:
            with open(path, "w") as f:
                f.write(str(exported))
        return exported

    def bridge_stats(self):
        stats = self.bridge.stats()
        return {
            "bridge_stats": stats,
            "api_diagnostics": self.api.diagnostics(),
        }

    def session_audit(self, session: APISession):
        is_open = session.is_open()
        session.checkpoint()
        return {
            "session_open": is_open,
            "checkpoint_taken": True,
        }

    def run_diagnostic_cycle(self, coordinate, proposition):
        session = self.api.open_session("diagnostics")
        request = APIRequest(
            request_id="diag-001",
            operation=OperationKind.INSPECT,
            coordinate=coordinate,
            proposition=proposition,
            budget=5000,
            deadline=10000,
            metadata={"purpose": "diagnostics"},
        )
        session.register_request(request)

        result = self.api.inspect(coordinate=coordinate)
        copilot_resp = self.api.copilot_query(
            f"Diagnose status of {coordinate}"
        )

        audit = self.session_audit(session)
        stats = self.bridge_stats()
        log = self.export_log()

        session.close()
        self.event_log.clear()

        return {
            "inspection": result,
            "copilot_diagnostic": copilot_resp,
            "audit": audit,
            "stats": stats,
            "log_size": len(log) if log else 0,
        }
\end{lstlisting}

% ============================================================
%  Section 8 — Formal Properties
% ============================================================
\section{Formal Properties}\label{sec:formal}

\begin{theorem}[Trust ceiling enforcement]\label{thm:trustceil}
For every API request processed by the $\Bridge$, the
resulting $\APIResp$ has trust $\tleq \Tcopilot$.
\end{theorem}
\begin{proof}
The bridge invokes the Copilot channel, whose responses
are capped at $\Tcopilot$ by the channel's
\texttt{apply\_trust\_ceiling} method.  The bridge then
calls its own \texttt{enforce\_trust\_ceiling} method,
which re-checks the bound.  The API response wraps the
channel response without elevating trust.  Therefore the
final trust is at most $\Tcopilot$.

In the corroboration case, the corroborating channel (e.g.,
solver) may produce evidence at $\Tsolver$.  However, this
evidence is attributed to the solver, not the bridge.  The
bridge's own output remains at $\Tcopilot$.
\end{proof}

\begin{theorem}[Corroboration soundness]\label{thm:corrobsound}
If corroboration succeeds, the corroborated evidence is
consistent with the Copilot suggestion and has trust at
least as high as the corroborating channel's ceiling.
\end{theorem}
\begin{proof}
Corroboration succeeds when the solver channel returns a
non-empty response for the same proposition.  By definition
of corroboration (Definition~\ref{def:corroboration}), the
responses are consistent.  The solver's response has trust
$\Tsolver > \Tcopilot$.  Therefore the corroborated evidence
has trust $\Tsolver$.

If corroboration fails (the solver returns empty or
contradictory evidence), the bridge returns the
uncorroborated Copilot response at $\Tcopilot$, which
is sound but lower-trust.
\end{proof}

\begin{theorem}[Event log completeness]\label{thm:logcomp}
Every bridge invocation generates a complete audit trail
in the $\APILog$.  Specifically, for each query:
\begin{enumerate}
\item A ``bridge\_query'' event is recorded.
\item If corroboration is attempted, a
      ``corroboration\_success'' or ``corroboration\_failed''
      event is recorded.
\item The final result is recorded.
\end{enumerate}
\end{theorem}
\begin{proof}
Inspection of the bridge implementation
(Listing~\ref{lst:bridge}) shows that
\texttt{event\_log.record} is called:
(1)~after the query (``bridge\_query'' event),
(2)~after corroboration (``corroboration\_success'' or
``corroboration\_failed''), and
(3)~after the result is produced (in the calling methods
\texttt{verify\_via\_bridge}, \texttt{construct\_via\_bridge},
etc.).  Therefore the log is complete.
\end{proof}

\begin{theorem}[Session isolation]\label{thm:sessiso}
Concurrent API sessions do not interfere.  Each session
maintains an independent request log, checkpoint state,
and lifecycle.
\end{theorem}
\begin{proof}
Each $\APISess$ instance has private fields: a request
list, a session id, an open/closed flag, and checkpoint
data.  No global mutable state is shared between sessions.
The $\APILog$ is shared but is append-only, so concurrent
appends do not cause data loss or corruption (under the
assumption of a thread-safe append operation).
\end{proof}

\begin{theorem}[Rate-limit compliance]\label{thm:ratelimit}
The bridge respects the global rate limit enforced by
$\APIRate$.  No bridge invocation exceeds the configured
capacity and refill rate.
\end{theorem}
\begin{proof}
All bridge operations go through $\JuGeoAPI$, which
applies the $\APIRate$ check before processing any request.
The rate limiter tracks tokens with a token-bucket algorithm:
each request consumes one token, and tokens refill at the
configured rate.  If no tokens are available, the request
is rejected.  Therefore the bridge cannot exceed the rate
limit.
\end{proof}

% ============================================================
%  Section 9 — Discussion
% ============================================================
\section{Discussion}\label{sec:discussion}

\paragraph{Performance overhead of corroboration.}
Corroboration doubles the cost of each Copilot query,
since a solver invocation is required.  In practice,
corroboration should be enabled selectively---for
high-stakes operations (e.g., \texttt{VERIFY},
\texttt{CONSTRUCT}) and disabled for low-stakes ones
(e.g., \texttt{INSPECT}).

\paragraph{Bridge as a bottleneck.}
The bridge serialises Copilot queries: each API request
that involves the Copilot channel goes through the bridge.
For high-throughput scenarios, the bridge can be replicated
(one per Copilot channel instance) and load-balanced.

\paragraph{Trust ceiling and corroboration tension.}
The bridge enforces $\Tcopilot$ on its own output but
allows corroboration to elevate trust via an independent
channel.  This design separates concerns: the bridge is
responsible for the Copilot layer, while the corroborating
channel is responsible for higher-trust evidence.

\paragraph{Event log volume.}
In a busy verification session, the event log can grow
large.  The \texttt{clear} method allows periodic pruning.
Exported logs can be persisted for offline analysis.

\paragraph{Extending the bridge.}
The bridge currently supports five operation kinds.  New
kinds (e.g., $\texttt{REFINE}$, $\texttt{ABSTRACT}$) can
be added by extending the prompt templates and registering
new operation kinds.

% ============================================================
%  Section 10 — Related Work
% ============================================================
\section{Related Work}\label{sec:related}

\paragraph{API gateway patterns.}
The bridge is analogous to an API
gateway~\cite{newman2015microservices} that mediates
between external clients and internal services.  Our
bridge adds trust enforcement, a concern absent in
traditional gateways.

\paragraph{Language model integration.}
OpenAI's function
calling~\cite{openai2023functions} and LangChain's tool
abstraction~\cite{chase2023langchain} provide similar
bridge functionality.  Our bridge differs in its formal
trust model and corroboration protocol.

\paragraph{Evidence composition.}
The federation protocol (Paper~88) composes evidence from
multiple channels.  The bridge integrates with federation
(Listing~\ref{lst:fedbridge}) to support the
\texttt{FEDERATE} operation kind.

\paragraph{Copilot advisors.}
Papers~80--84 and~99 of this series use the bridge
(directly or indirectly) for Copilot queries.  This paper
provides the formal foundation for those uses.

\paragraph{Trust calibration.}
Ribeiro \etal{}~\cite{ribeiro2020trust} study trust in
AI predictions.  Our trust lattice and ceiling mechanism
provide a formal counterpart.

\paragraph{API design for verification.}
Boogie~\cite{barnett2006boogie} and Why3~\cite{filliatre2013why3}
expose verification APIs.  Our API layer adds trust
annotations and Copilot integration.

% ============================================================
%  Section 11 — Conclusion
% ============================================================
\section{Conclusion}\label{sec:conclusion}

We have provided a comprehensive formal treatment of the
$\Bridge$, the central integration point between the
Copilot channel and the \JuGeo{} API layer.  The bridge:
\begin{itemize}
\item Translates API operations into Copilot queries.
\item Enforces the $\Tcopilot$ trust ceiling on all outputs.
\item Optionally corroborates Copilot suggestions via
      higher-trust channels.
\item Logs all interactions for auditing.
\item Respects session isolation and rate limits.
\end{itemize}
Five theorems establish trust ceiling enforcement,
corroboration soundness, event log completeness, session
isolation, and rate-limit compliance.  The bridge is the
foundation upon which all Copilot advisors in the \JuGeo{}
series are built.

This paper concludes the first hundred papers in the
Judgment Geometry series.  Future work includes extending
the bridge to support streaming responses, adding
fine-grained trust annotations per proposition, and
integrating with external knowledge graphs for richer
corroboration.

% ============================================================
%  Bibliography
% ============================================================
\begin{thebibliography}{25}

\bibitem{newman2015microservices}
S.~Newman,
\emph{Building Microservices}.
O'Reilly, 2015.

\bibitem{openai2023functions}
OpenAI,
``Function calling and other API updates,''
\emph{OpenAI Blog}, 2023.

\bibitem{chase2023langchain}
H.~Chase,
``LangChain: building applications with LLMs through
composability,''
2023.

\bibitem{ribeiro2020trust}
M.~T.~Ribeiro, S.~Singh, and C.~Guestrin,
``'Why should I trust you?': explaining the predictions
of any classifier,''
\emph{KDD}, 2016.

\bibitem{barnett2006boogie}
M.~Barnett, B.-Y.~E.~Chang, R.~DeLine, B.~Jacobs, and
K.~R.~M.~Leino,
``Boogie: a modular reusable verifier for object-oriented
programs,''
\emph{FMCO}, 2006.

\bibitem{filliatre2013why3}
J.-C.~Filli\^{a}tre and A.~Paskevich,
``Why3 --- where programs meet provers,''
\emph{ESOP}, 2013.

\bibitem{young2024jugeo}
H.~Young,
``Judgment geometry: a framework for trust-calibrated
verification,''
\emph{JuGeo Technical Reports}, Paper~1, 2024.

\bibitem{young2024channels}
H.~Young,
``Evidence channels in judgment geometry,''
\emph{JuGeo Technical Reports}, Paper~88, 2024.

\bibitem{young2024advisors}
H.~Young,
``Copilot advisors for heap, scope, import, and callable
analysis,''
\emph{JuGeo Technical Reports}, Papers~80--81, 2024.

\bibitem{young2024research}
H.~Young,
``Copilot research and experiment advisors,''
\emph{JuGeo Technical Reports}, Paper~82, 2024.

\bibitem{young2024optimization}
H.~Young,
``Copilot optimization advisor,''
\emph{JuGeo Technical Reports}, Paper~83, 2024.

\bibitem{young2024lowering}
H.~Young,
``Copilot lowering hints and IR advisors,''
\emph{JuGeo Technical Reports}, Paper~84, 2024.

\bibitem{young2024connection}
H.~Young,
``Copilot connection hooks and health checks,''
\emph{JuGeo Technical Reports}, Paper~99, 2024.

\bibitem{young2024lifecycle}
H.~Young,
``Kernel lifecycle and health monitoring in judgment geometry,''
\emph{JuGeo Technical Reports}, Paper~93, 2024.

\bibitem{demoura2008z3}
L.~de~Moura and N.~Bj{\o}rner,
``Z3: an efficient SMT solver,''
\emph{TACAS}, 2008.

\bibitem{moura2021lean4}
L.~de~Moura and S.~Ullrich,
``The Lean~4 theorem prover and programming language,''
\emph{CADE}, 2021.

\bibitem{polu2020gptf}
S.~Polu and I.~Sutskever,
``Generative language modeling for automated theorem proving,''
\emph{arXiv preprint arXiv:2009.03393}, 2020.

\bibitem{lample2022hypertree}
G.~Lample \etal{},
``HyperTree proof search for neural theorem proving,''
\emph{NeurIPS}, 2022.

\bibitem{bansal2019holist}
K.~Bansal, S.~Loos, M.~Rabe, C.~Szegedy, and S.~Wilcox,
``HOList: an environment for machine learning of higher-order
logic theorem proving,''
\emph{ICML}, 2019.

\bibitem{yang2019learning}
K.~Yang and J.~Deng,
``Learning to prove theorems via interacting with proof
assistants,''
\emph{ICML}, 2019.

\bibitem{first2022diversity}
E.~First, M.~Rabe, T.~Ringer, and Y.~Brun,
``Diversity-driven automated formal verification,''
\emph{ICSE}, 2022.

\bibitem{jiang2022thor}
A.~Q.~Jiang, W.~Li, S.~Tworber, and Y.~Wu,
``Thor: wielding hammers to integrate language models and
automated theorem provers,''
\emph{NeurIPS}, 2022.

\end{thebibliography}

% ============================================================
%  Appendix
% ============================================================
\appendix

\section{Auxiliary Definitions and Proofs}\label{app:proofs}

\begin{lemma}[Bridge natural transformation well-definedness]%
\label{lem:bridgenatwd}
The mapping $\eta_o \colon F_A(o) \to F_C(o)$ is
well-defined for every operation kind~$o$.
\end{lemma}
\begin{proof}
For each $o \in \OpKind$, the bridge constructs a prompt
template specific to~$o$ and produces an $\EvReq$ with
the appropriate \texttt{proposition\_kind}.  The mapping
is total (defined for all valid $\APIReq$ values) and
functional (each request maps to a unique evidence request).
\end{proof}

\begin{lemma}[Naturality of the bridge]%
\label{lem:bridgenaturality}
For every morphism $o \to o'$ in $\mathbf{API}$,
$\eta_{o'} \circ F_A(o \to o') = F_C(o \to o') \circ \eta_o$.
\end{lemma}
\begin{proof}
Let $r \in F_A(o)$ be an API request.  Restricting $r$ to
kind~$o'$ gives $r' = F_A(o \to o')(r)$.  Translating $r'$
gives $\eta_{o'}(r')$, an evidence request for kind~$o'$.

Alternatively, translating $r$ gives $\eta_o(r)$, an
evidence request for kind~$o$.  Restricting to kind~$o'$
gives $F_C(o \to o')(\eta_o(r))$.  The bridge's prompt
construction ensures that restriction commutes with
translation: both paths produce the same evidence request
for the sub-operation~$o'$.
\end{proof}

\begin{lemma}[Trust ceiling is a natural bound]%
\label{lem:trustnatbound}
The trust ceiling $\Tcopilot$ is uniform across all
operation kinds.
\end{lemma}
\begin{proof}
The ceiling is determined by the Copilot channel
configuration, which is independent of the operation kind.
The bridge applies the same ceiling regardless of whether
the operation is VERIFY, CONSTRUCT, DESCEND, FEDERATE, or
INSPECT.
\end{proof}

\begin{lemma}[Event log append-only property]%
\label{lem:logappend}
The $\APILog$ is append-only: events can be recorded and
exported, but not modified after recording.
\end{lemma}
\begin{proof}
The \texttt{record} method appends to an internal list.
The \texttt{export} method returns a copy.  The
\texttt{clear} method resets the list but does not modify
existing exported copies.  No method modifies previously
recorded events.
\end{proof}

\begin{lemma}[Rate limiter fairness]\label{lem:ratefair}
The token-bucket rate limiter treats all callers equally:
each request consumes one token regardless of the operation
kind or caller identity.
\end{lemma}
\begin{proof}
The $\APIRate$ class maintains a single token bucket
parameterised by capacity and refill rate.  The
\texttt{consume} method decrements the bucket by one for
each request.  No caller or operation receives preferential
treatment.
\end{proof}

\begin{lemma}[Corroboration does not violate trust ceiling]%
\label{lem:corrobnoceil}
Even when corroboration introduces solver evidence at
$\Tsolver$, the bridge's own output remains at $\Tcopilot$.
\end{lemma}
\begin{proof}
The bridge attributes corroborated evidence to the solver
channel, not to itself.  The bridge's output (the hint or
suggestion) remains at $\Tcopilot$.  The corroborated
evidence is a separate artifact with its own trust
annotation.
\end{proof}

\begin{corollary}[Bridge trust ceiling is inviolable]%
\label{cor:bridgeinviolable}
Under no configuration of the bridge does the bridge's
own output exceed $\Tcopilot$.
\end{corollary}
\begin{proof}
By Lemma~\ref{lem:bridgeceil} (basic trust ceiling),
Lemma~\ref{lem:clampidemp} (clamping is idempotent),
and Lemma~\ref{lem:corrobnoceil} (corroboration does
not violate ceiling), the bridge's output trust is at most
$\Tcopilot$ in all cases.
\end{proof}

\begin{lemma}[Session checkpoint preserves state]%
\label{lem:sesschkpt}
After calling \texttt{session.checkpoint()}, the session's
request list and metadata are serialised.  A subsequent
\texttt{close()} does not lose the checkpointed data.
\end{lemma}
\begin{proof}
The \texttt{checkpoint} method serialises the current state
to persistent storage.  The \texttt{close} method marks
the session as closed but does not delete the checkpointed
data.  Therefore the checkpoint is preserved.
\end{proof}

\begin{proposition}[Bridge correctness summary]%
\label{prop:bridgecorrect}
The $\Bridge$ correctly mediates between the API layer and
the channel layer.  Specifically:
\begin{enumerate}
\item Every API request that involves Copilot is translated
      to an evidence request (Lemma~\ref{lem:bridgenatwd}).
\item The translation is coherent across operation kinds
      (Lemma~\ref{lem:bridgenaturality}).
\item Trust is bounded by $\Tcopilot$
      (Theorem~\ref{thm:trustceil}).
\item Corroboration, when successful, provides higher-trust
      evidence (Theorem~\ref{thm:corrobsound}).
\item All interactions are logged
      (Theorem~\ref{thm:logcomp}).
\item Sessions are isolated
      (Theorem~\ref{thm:sessiso}).
\item Rate limits are respected
      (Theorem~\ref{thm:ratelimit}).
\end{enumerate}
\end{proposition}
\begin{proof}
Each item is established by the referenced theorem or lemma.
Together, they show that the bridge is a correct, trust-safe,
auditable, and rate-limited integration layer.
\end{proof}

\begin{remark}[Bridge as the capstone]
Paper~100 concludes the first century of the Judgment Geometry
series.  The $\Bridge$ is the keystone component that enables
all Copilot-assisted verification: without the bridge, the
language model's suggestions cannot enter the trust-calibrated
verification pipeline.  With the bridge, those suggestions are
safely bounded, optionally corroborated, fully logged, and
rate-limited.
\end{remark}

\end{document}
